\documentclass[11pt,a4paper]{article}
\usepackage[margin=2cm]{geometry}
\usepackage{amsmath,amssymb}
\usepackage{hyperref}
\usepackage{enumitem}
\usepackage{booktabs}
\usepackage{graphicx}
\usepackage[T1]{fontenc}
\usepackage{lmodern}
\usepackage{microtype}

\setlength{\parskip}{0.4em}
\setlength{\parindent}{0pt}

\title{\vspace{-2em}\textbf{QDay Prize Submission Brief}\\[0.3em]
\large End-to-End Compilable Quantum Elliptic Curve Discrete Logarithm}
\author{Diego Polimeni\\\small\texttt{diego.polimeni.work@gmail.com}}
\date{}

\begin{document}
\maketitle
\vspace{-1.5em}

\section*{Approach}

We implement \textbf{Shor's algorithm for the Elliptic Curve Discrete Logarithm Problem (ECDLP)} using \textbf{Qrisp}, a high-level quantum programming language built on top of Qiskit. Given a generator point $G$ and public key $P = kG$ on an elliptic curve $y^2 = x^3 + 7 \pmod{p}$, the algorithm recovers the private key $k$ via quantum phase estimation on the group operation.

The core algorithm prepares two $n$-qubit QPE registers in superposition, performs $2n$ controlled EC point additions (multiplying by both $G$ and $P$), applies inverse QFT, and extracts $k$ from the measurement outcomes via continued-fraction expansion. This follows the standard Shor ECDLP construction described in our paper~\cite{polimeni2025}.

\section*{Techniques and Design Choices}

\textbf{Static circuit compilation.}  
As detailed in~\cite{polimeni2025}, we realize all EC arithmetic---point addition, point doubling, Kaliski modular inversion, Montgomery modular multiplication---as compilable quantum circuits using Qrisp's \texttt{QuantumModulus} type and in-place modular arithmetic primitives. Every operation is decomposed into elementary gates (H, CNOT, T, T$^\dagger$) at compile time.

\textbf{Dynamic compilation via JAX (\texttt{jasp}).}  
Beyond the static approach, we introduce \textbf{JAX-traceable dynamic compilation} through Qrisp's \texttt{jasp} subsystem. The quantum algorithm is traced through \texttt{jax.jit}, enabling dynamic circuit generation, loop unrolling, and optimization. This is critical for scaling to large bit sizes where static unrolling would be intractable.

\textbf{BigInteger support.}  
To handle prime fields exceeding the native 64-bit JAX integer limit, we developed a \texttt{BigInteger} class---a fixed-width array of 32-bit limbs registered as a JAX pytree node. This enables compilation for arbitrarily large primes (up to 256-bit and beyond) while remaining fully JAX-traceable. All classical computation uses Python \texttt{int} values; \texttt{BigInteger} values appear only in \texttt{QuantumModulus} register construction and internal operator dispatch.

\textbf{Scalable resource estimation.}  
Every controlled EC addition in the Shor loop has \emph{identical gate structure} regardless of the classical point---only the modulus bit-width matters. We exploit this by counting gates for a \emph{single} controlled addition via Qrisp's \texttt{@count\_ops} decorator and computing:
\[
\text{Total gates} = 2n \times \text{(single EC addition)} + \text{overhead}
\]
where overhead covers Hadamard preparation, initial encoding, inverse QFT, and measurement. This makes 256-bit resource estimation feasible on standard hardware without tracing 512 full EC additions.

To make the repository structure explicit, we separate the two resource-estimation use cases into two notebooks: \texttt{resource\_estimation\_all\_curves.ipynb} compiles all QDay prize curves (4--21 bit), while \texttt{EC\_discrete\_log\_compilation.ipynb} focuses on larger generated curves and the scaling path up to 256-bit.

\section*{Requirements and Execution}

\begin{itemize}[nosep,leftmargin=1.2em]
    \item \textbf{Software}: Python 3.10+, Qrisp (development branch with PRs \href{https://github.com/eclipse-qrisp/qrisp/pull/495}{\#495} and \href{https://github.com/eclipse-qrisp/qrisp/pull/505}{\#505}, both pending merge), JAX, NumPy, SymPy, Matplotlib.
    \item \textbf{Hardware}: Standard laptop/desktop. Resource estimation for 256-bit takes minutes; full simulation is limited to small curves ($\leq$5-bit) due to exponential state-space scaling.
    \item \textbf{Execution}: Install via \texttt{pip install -r requirements.txt}. Run \texttt{resource\_estimation\_all\_curves.ipynb} for the full QDay prize curve set, and run \texttt{EC\_discrete\_log\_compilation.ipynb} for compilation and resource estimation on larger curves up to 256-bit.
    \item \textbf{Validation}: Unit tests (\texttt{pytest tests/ -v}) verify correctness of all quantum primitives against classical reference implementations using Qrisp's boolean simulator.
\end{itemize}

\section*{Potential Drawbacks}

\begin{itemize}[nosep,leftmargin=1.2em]
    \item \textbf{No fault-tolerant hardware execution}: Current quantum hardware cannot execute circuits of this depth. The submission demonstrates correctness and compilation feasibility, not hardware execution.
    \item \textbf{No T-gate optimization}: The current compilation does not apply T-gate reduction techniques (e.g., T-factory scheduling, magic state distillation budgeting). Gate counts represent unoptimized decompositions.
    \item \textbf{Qrisp dependency on development branches}: The implementation requires two PRs not yet merged into Qrisp's main branch.
    \item \textbf{Static point encoding}: The current approach encodes classical EC points statically. Windowed or semi-classical QFT variants could reduce qubit count but are not yet implemented.
\end{itemize}

\begin{thebibliography}{1}
\bibitem{polimeni2025} D.~Polimeni, R.~Seidel. \emph{End-to-end compilable implementation of quantum elliptic curve logarithm in Qrisp}. arXiv:2501.10228 (2025).
\end{thebibliography}

\end{document}
